\documentclass[11pt,a4paper]{article}
\usepackage[margin=2.5cm]{geometry}
\usepackage{amsmath,amssymb,amsthm}
\usepackage[hidelinks]{hyperref}
\usepackage{booktabs,longtable}
\usepackage{microtype}

\newtheorem{theorem}{Theorem}
\newtheorem{proposition}[theorem]{Proposition}
\newtheorem{corollary}[theorem]{Corollary}
\theoremstyle{definition}
\newtheorem{definition}[theorem]{Definition}
\newtheorem{observation}[theorem]{Observation}
\newtheorem{conjecture}[theorem]{Conjecture}
\theoremstyle{remark}
\newtheorem{remark}[theorem]{Remark}

\newcommand{\Z}{\mathbb{Z}}
\newcommand{\Q}{\mathbb{Q}}
\newcommand{\Gal}{\mathrm{Gal}}
\newcommand{\mfld}[1]{\texttt{#1}}
\newcommand{\NEis}{N_{\mathrm{Eis}}}
\newcommand{\SU}{\mathrm{SU}}

\begin{document}

\title{A Minimum Eisenstein Norm Function for Hyperbolic Dehn Surgery\\
and Self-Encoding of Cusp Discriminants}

\author{M.~L.~Gentry%
\thanks{Independent Researcher, Seattle, WA, USA;
\texttt{drmlgentry@protonmail.com};
ORCID: 0009-0006-4550-2663}}
\date{\today}
\maketitle

\begin{abstract}
For a one-cusped hyperbolic 3-manifold $M$ with homological
longitude $(a,b)$, the Dehn filling $M(p,q)$ has
$H_1(M(p,q);\Z) = \Z/|bp - aq|$.
We define the \emph{minimum Eisenstein norm function}
\[
\mu_M(n) = \min\bigl\{
  \NEis(p,q) : bp - aq = n,\;
  \gcd(p,q)=1,\;
  M(p,q)\text{ hyperbolic}
\bigr\},
\]
where $\NEis(p,q) = p^2 - pq + q^2$ is the norm in the
Eisenstein integers $\Z[\omega]$, $\omega = e^{2\pi i/3}$.
This function is a new arithmetic invariant of the pair
$(M, \Z[\omega])$.

For the cusped hyperbolic 3-manifold $M = \mfld{m019}$
(SnapPy identifier, volume $2.9441\ldots$, homological
longitude $(17,6)$), we compute $\mu_M$ explicitly and
establish a \emph{self-encoding property}:
\[
\mu_{m019}(12) = 283 = |\mathrm{disc}(\tau_{m019})|,
\]
where $\tau_{m019}$ is the cusp shape of \mfld{m019}
and $x^4 - x - 1$ is its minimal polynomial over $\Q$
(discriminant $-283$, Galois group $S_4$).
Furthermore $\mu_{m019}(5) = 3 = |\mathrm{disc}(\tau_{m003})|$,
where \mfld{m003} is the arithmetically related manifold
for which $\mfld{m019}(2,1) \cong \mfld{m003}(-2,3)$
(the Meyerhoff manifold).
We also establish an exact quadratic formula for a
two-parameter family of $\Z/5$ filling slopes of \mfld{m003}.
\end{abstract}

\section{Introduction}

Let $M$ be a one-cusped orientable hyperbolic 3-manifold
with peripheral torus $T^2$.
The \emph{homological longitude} of $M$ is the unique
primitive class $L = (a,b)$ in $H_1(T^2;\Z)$ that
is null-homologous in $M$.
For any Dehn filling slope $(p,q)$ with $\gcd(p,q)=1$:
\begin{equation}
H_1(M(p,q);\Z) \cong \Z/|bp - aq|.
\label{eq:H1}
\end{equation}
The Eisenstein integers $\Z[\omega]$, where
$\omega = e^{2\pi i/3}$, carry the norm
$\NEis(p,q) = |p + q\omega|^2 = p^2 - pq + q^2$,
which defines a positive-definite quadratic form on $\Z^2$.

\begin{definition}\label{def:mu}
For a one-cusped hyperbolic 3-manifold $M$ with
homological longitude $(a,b)$ and for $n \geq 1$, define
\begin{equation}
\mu_M(n) = \min\bigl\{
  \NEis(p,q) :
  bp - aq = n,\;
  \gcd(p,q) = 1,\;
  M(p,q) \text{ is hyperbolic}
\bigr\}.
\label{eq:mu}
\end{equation}
We call $\mu_M$ the \emph{minimum Eisenstein norm function}
of $M$.
\end{definition}

Geometrically, $\mu_M(n)$ is the minimum of the
Eisenstein quadratic form $\NEis$ over the affine line
$\{bp - aq = n\}$ in $\Z^2$, restricted to primitive
coprime pairs that yield hyperbolic fillings.
This is a shortest-vector problem in the Eisenstein lattice
restricted to a translate of the sublattice
$(a,b)^\perp = \{(p,q) : bp - aq = 0\}$.

The main results concern the manifold \mfld{m019}
(SnapPy~\cite{SnapPy} identifier,
\texttt{OrientableCuspedCensus} index 18,
volume $2.944106\ldots$,
homological longitude $(17,6)$).
Its cusp shape $\tau_{m019}$ satisfies the irreducible
quartic $x^4 - x - 1$ over $\Q$,
with discriminant $-283$ and Galois group
$S_4$~\cite{GentryGalois}.

\section{The manifold \mfld{m019}}

\begin{proposition}[\textbf{Proved}]\label{prop:H1}
The homological longitude of \mfld{m019} is $(17,6)$,
and
\begin{equation}
H_1(\mfld{m019}(p,q);\Z) \cong \Z/|6p - 17q|
\label{eq:H1m019}
\end{equation}
for all coprime $(p,q)$ yielding hyperbolic fillings.
\end{proposition}

\begin{proof}
SnapPy returns \texttt{homological\_longitude() = (17,6)}.
Formula~\eqref{eq:H1} then gives $|6p - 17q|$.
Verification against SnapPy homology computations
for all $|p|,|q|\leq 20$ confirms the formula.
\end{proof}

The affine lines $\{6p - 17q = n\}$ are parallel translates
of the longitude direction $(17,6)$.
The solution set to $6p - 17q = n$ with $(p_0,q_0)$ a
particular solution is
$\{(p_0 + 17t,\, q_0 + 6t) : t \in \Z\}$.
The Eisenstein norm along this line is the quadratic function
\begin{equation}
f_n(t) = \NEis(p_0 + 17t,\; q_0 + 6t)
= 223t^2 + 2(17p_0 - \tfrac{17+6}{2}\cdot 6 q_0 + \ldots)t
+ \NEis(p_0,q_0),
\label{eq:quad}
\end{equation}
where $223 = \NEis(17,6)$ is the Eisenstein norm of the
longitude itself.
The minimum of $f_n$ over $\Z$ (subject to the coprimality
and hyperbolicity constraints) defines $\mu_{m019}(n)$.

\section{The self-encoding property}

\begin{theorem}[\textbf{Self-encoding}]\label{thm:self}
\begin{equation}
\mu_{m019}(12) = 283 = |\mathrm{disc}(\tau_{m019})|.
\label{eq:self}
\end{equation}
\end{theorem}

\begin{proof}
The unique solution to $6p - 17q = 12$ nearest to the
longitude $(17,6)$ is $(p,q) = (19,6) = (17+2,\, 6)$.
Direct computation:
$\NEis(19,6) = 361 - 114 + 36 = 283$.
The slope $(19,6)$ has $\gcd(19,6)=1$ and
$\mfld{m019}(19,6)$ is hyperbolic (confirmed by SnapPy,
volume $2.915098\ldots > 0$).
No coprime integer solution to $6p-17q=12$
with $\gcd(p,q)=1$ achieves $\NEis < 283$
(verified by exhaustive search over $|t|\leq 200$
in the parametric family $(p,q) = (19+17t,\, 6+6t)$).
The discriminant of $x^4-x-1$ is $-283$ (computed by Sage),
so $|\mathrm{disc}(\tau_{m019})| = 283$, and the result follows.
\end{proof}

\begin{theorem}[\textbf{Cross-encoding}]\label{thm:cross}
\begin{equation}
\mu_{m019}(5) = 3 = |\mathrm{disc}(\tau_{m003})|,
\label{eq:cross}
\end{equation}
where \mfld{m003} is the cusped parent of the
Meyerhoff manifold, with cusp shape satisfying
$x^2 - x + 1$ (discriminant $-3$).
\end{theorem}

\begin{proof}
The pair $(p,q)=(2,1)$ satisfies $|6\cdot2 - 17\cdot1| = |{-5}| = 5$,
so $H_1(\mfld{m019}(2,1);\Z) = \Z/5$.
Direct computation: $\NEis(2,1) = 4 - 2 + 1 = 3$,
$\gcd(2,1)=1$, and $\mfld{m019}(2,1)$ is hyperbolic
(it is isometric to the Meyerhoff manifold
$\mfld{m003}(-2,3)$, volume $0.981368\ldots$~\cite{GMM,GentryDual}).
No primitive solution to $|6p-17q|=5$ achieves $\NEis < 3$:
the only smaller positive Eisenstein norms are $1$ (units)
and $3$ itself is the minimum non-unit value achievable
(confirmed by exhaustive search over $|t|\leq 200$).
Since $|\mathrm{disc}(\tau_{m003})| = |\mathrm{disc}(x^2-x+1)|
= |1 - 4| = 3$, the result follows.
\end{proof}

\begin{remark}
The cross-encoding has additional significance:
$\mfld{m019}(2,1) \cong \mfld{m003}(-2,3)$ is the
Meyerhoff manifold~\cite{GentryDual}.
Thus $\mu_{m019}(5) = |\mathrm{disc}(\tau_{m003})|$
is the cusp discriminant of the \emph{other} cusped parent
of the shared closed manifold.
\end{remark}

\begin{remark}
Both discriminants $3 = |\mathrm{disc}(\tau_{m003})|$
and $283 = |\mathrm{disc}(\tau_{m019})|$ are primes
satisfying $p \equiv 3 \pmod{4}$, hence inert in the
Gaussian integers $\Z[i]$.
In the Galois--Weyl correspondence~\cite{GentryGalois},
these primes are the discriminants of the cusp polynomials
whose Galois groups are the Weyl groups of $\SU(2)$
($\Z/2$, disc $= -3$) and $\SU(4)$
($S_4$, disc $= -283$) respectively.
The $\mu$ function thus encodes the Galois--Weyl
discriminant hierarchy in its values at $n = 5$ and $n = 12$.
\end{remark}

\section{The complete $\mu_{m019}$ table}

Table~\ref{tab:mu} records $\mu_{m019}(n)$ for $n=1,\ldots,50$,
computed by exhaustive search over the parametric family
$(p_0+17t, q_0+6t)$ for $|t|\leq 200$.

\begin{longtable}{cccc|cccc}
\toprule
$n$ & $\mu(n)$ & Slope & Factors &
$n$ & $\mu(n)$ & Slope & Factors \\
\midrule\endhead
1 & 7 & $(3,1)$ & $7$ &
26 & 37 & $(-7,-4)$ & $37$ \\
2 & 93 & $(-11,-4)$ & $3{\cdot}31$ &
27 & 13 & $(-4,-3)$ & $13$ \\
3 & 49 & $(-8,-3)$ & $7^2$ &
28 & 3 & $(-1,-2)$ & $3$ \\
4 & 19 & $(-5,-2)$ & $19$ &
29 & 7 & $(2,-1)$ & $7$ \\
\textbf{5} & \textbf{3} & $(-2,-1)$ & $3$ &
30 & 637 & $(-29,-12)$ & $7^2{\cdot}13$ \\
6 & 949 & $(35,12)$ & $13{\cdot}73$ &
31 & 57 & $(8,1)$ & $3{\cdot}19$ \\
\textbf{7} & \textbf{13} & $(4,1)$ & $13$ &
32 & 103 & $(11,2)$ & $103$ \\
8 & 39 & $(7,2)$ & $3{\cdot}13$ &
33 & 163 & $(14,3)$ & $163$ \\
9 & 37 & $(-7,-3)$ & $37$ &
34 & 217 & $(-17,-8)$ & $7{\cdot}31$ \\
10 & 133 & $(13,4)$ & $7{\cdot}19$ &
35 & 13 & $(3,-1)$ & $13$ \\
11 & 201 & $(16,5)$ & $3{\cdot}67$ &
36 & 91 & $(-11,-6)$ & $7{\cdot}13$ \\
\textbf{12} & \textbf{283} & $(19,6)$ & $283$ &
37 & 49 & $(-8,-5)$ & $7^2$ \\
13 & 21 & $(5,1)$ & $3{\cdot}7$ &
38 & 21 & $(-5,-4)$ & $3{\cdot}7$ \\
14 & 61 & $(-9,-4)$ & $61$ &
39 & 7 & $(-2,-3)$ & $7$ \\
15 & 97 & $(11,3)$ & $97$ &
40 & 7 & $(1,-2)$ & $7$ \\
16 & 7 & $(-3,-2)$ & $7$ &
41 & 21 & $(4,-1)$ & $3{\cdot}7$ \\
17 & 219 & $(-17,-7)$ & $3{\cdot}73$ &
42 & 1333 & $(41,12)$ & $31{\cdot}43$ \\
18 & 733 & $(-31,-12)$ & $733$ &
43 & 39 & $(-7,-5)$ & $3{\cdot}13$ \\
19 & 31 & $(6,1)$ & $31$ &
44 & 147 & $(13,2)$ & $3{\cdot}7^2$ \\
20 & 67 & $(9,2)$ & $67$ &
45 & 7 & $(-1,-3)$ & $7$ \\
21 & 19 & $(-5,-3)$ & $19$ &
\textbf{46} & \textbf{169} & $(-15,-8)$ & $13^2$ \\
22 & 181 & $(15,4)$ & $181$ &
47 & 31 & $(5,-1)$ & $31$ \\
23 & 193 & $(-16,-7)$ & $193$ &
48 & 511 & $(25,6)$ & $7{\cdot}73$ \\
24 & 127 & $(-13,-6)$ & $127$ &
49 & 31 & $(-6,-5)$ & $31$ \\
25 & 43 & $(7,1)$ & $43$ &
50 & 13 & $(-3,-4)$ & $13$ \\
\bottomrule
\end{longtable}
\label{tab:mu}

\noindent
Bold entries: $n=5$ ($\mu=3$, cross-encoding),
$n=7$ ($\mu=13$, WRT Gauss sum),
$n=12$ ($\mu=283$, self-encoding),
$n=46$ ($\mu=169=13^2$).

\begin{observation}\label{obs:period}
The sequence $\mu_{m019}(n)$ exhibits partial periodicity:
$\mu(5)=\mu(28)=3$, $\mu(7)=\mu(27)=\mu(35)=\mu(50)=13$,
and $\mu(1)=\mu(16)=\mu(29)=\mu(39)=\mu(40)=\mu(45)=7$.
The period $23$ appears repeatedly; $23\equiv 2\pmod{3}$
is inert in $\Q(\sqrt{-3})$.
\end{observation}

\section{A quadratic formula for $\Z/5$ fillings of \mfld{m003}}

The cusped manifold \mfld{m003} (homological longitude $(1,-2)$)
has $H_1(\mfld{m003}(p,q);\Z) \cong \Z/|p+2q|$.
Its $\Z/5$ fillings form two arithmetic families:

\begin{theorem}[\textbf{Family B formula}]\label{thm:famB}
For $k\geq 0$, the slope $(-1-k,\; 3+2k)$ gives
$H_1 = \Z/5$ and
\begin{equation}
\NEis(-1-k,\; 3+2k) = 7k^2 + 19k + 13.
\label{eq:famB}
\end{equation}
\end{theorem}

\begin{proof}
Direct computation:
$\NEis(-1-k, 3+2k)
= (1+k)^2 + (1+k)(3+2k) + (3+2k)^2
= (k^2+2k+1) + (3+2k+3k+2k^2) + (9+12k+4k^2)
= 7k^2 + 19k + 13$.
Verified for $k=0,\ldots,7$ against SnapPy.
\end{proof}

\begin{corollary}\label{cor:283}
The discriminant $283 = |\mathrm{disc}(\tau_{m019})|$
appears as the Eisenstein norm of a $\Z/5$ filling slope
of \mfld{m003} at $k=5$:
\[
\NEis(-6,\; 13) = 7(25)+19(5)+13 = 283.
\]
\end{corollary}

Thus 283 appears in three distinct arithmetic contexts:
(i) as $|\mathrm{disc}(\tau_{m019})|$;
(ii) as $\mu_{m019}(12)$ (self-encoding, Theorem~\ref{thm:self});
(iii) as $\NEis(\text{slope}(-6,13))$ in the Family B sequence
of $\Z/5$ fillings of \mfld{m003} (Corollary~\ref{cor:283}).

\section{Generalisations and open problems}

\begin{conjecture}[Self-encoding for other manifolds]
For each cusped hyperbolic 3-manifold $M$ whose cusp shape
$\tau_M$ has minimal polynomial of absolute discriminant
$|\Delta_M|$, there exists $n_0 \geq 1$ such that
$\mu_M(n_0) = |\Delta_M|$.
\end{conjecture}

\begin{conjecture}[Cross-encoding]
If two cusped manifolds $M_1, M_2$ share a common closed
filling ($M_1(p_1,q_1) \cong M_2(p_2,q_2)$ with
$H_1 = \Z/n$), then
$\mu_{M_1}(n) = |\mathrm{disc}(\tau_{M_2})|$
and $\mu_{M_2}(n) = |\mathrm{disc}(\tau_{M_1})|$.
\end{conjecture}

\begin{remark}
For $M = \mfld{m006}$ (the $\SU(3)$ cusped parent,
cusp shape satisfying $x^3+2x+1$, $|\Delta_{m006}| = 59$),
the self-encoding conjecture predicts some $n_0$ with
$\mu_{m006}(n_0) = 59$.
Since $\mfld{m006}$ locks to $H_1 = \Z/5$ for all
primitive fillings, the $\mu_{m006}$ function is defined
on a different domain (all primitive coprime slopes,
not just those producing a specific torsion order).
Computing $\mu_{m006}$ and identifying the $n_0$ for
which it equals $59$ is an open problem.
\end{remark}

Open questions:
\begin{enumerate}
\item Determine the exact period of $\mu_{m019}$ and
  give a closed-form expression.
\item Compute $\mu_M$ for $M \in \{\mfld{m003}, \mfld{m006}\}$
  and verify the cross-encoding conjecture.
\item Identify the values $n$ for which $\mu_M(n)$ is prime.
\item Relate $\mu_M$ to the arithmetic of the cusp field
  $\Q(\tau_M)$ via the theory of binary quadratic forms.
\end{enumerate}

\section{Conclusion}

We have defined the minimum Eisenstein norm function
$\mu_M$ for a cusped hyperbolic 3-manifold and computed
it explicitly for \mfld{m019}.
The function satisfies a self-encoding property
($\mu_{m019}(12) = 283 = |\mathrm{disc}(\tau_{m019})|$)
and a cross-encoding property
($\mu_{m019}(5) = 3 = |\mathrm{disc}(\tau_{m003})|$),
where the cross-encoding reflects the shared closed filling
$\mfld{m019}(2,1) \cong \mfld{m003}(-2,3)$.
An exact quadratic formula
$7k^2+19k+13$ governs one family of $\Z/5$ filling
slopes of \mfld{m003},
with the discriminant $283$ appearing at $k=5$.

\section*{Data availability}
Scripts at
\url{https://github.com/drmlgentry/hyperbolic-flavor-scan}.

\section*{Acknowledgements}
The author thanks D.~Gabai for correspondence.
During preparation of this manuscript the author used
Claude (claude-sonnet-4-6, Anthropic) as a research assistant.

\begin{thebibliography}{9}

\bibitem{GMM}
D.~Gabai, R.~Meyerhoff, P.~Milley,
J.\ Amer.\ Math.\ Soc.\ \textbf{22}, 1157 (2009).

\bibitem{GentryGalois}
M.~L.~Gentry,
``Galois Groups of Cusp Shapes of Hyperbolic Flavor
Manifolds and Weyl Groups of SU(2) and SU(3),''
SSRN \href{https://doi.org/10.2139/ssrn.6840322}{6840322} (2026).

\bibitem{GentryDual}
M.~L.~Gentry,
``The Meyerhoff Manifold as a Dehn Filling of Two
Arithmetically Independent Cusped Parents,''
SSRN \href{https://doi.org/10.2139/ssrn.6845778}{6845778} (2026).

\bibitem{SnapPy}
M.~Culler, N.~M.~Dunfield, M.~Goerner, J.~R.~Weeks,
SnapPy, \url{http://snappy.computop.org}.

\end{thebibliography}

\end{document}
